Phase II design is deliberately out of scope. What Phase I supports is
a single directional statement:

\begin{quote}
The open capability is \emph{multiple independently retrievable
persistent representations under sequential experience}. Future work
should investigate architectures that make sequential, late-arriving
patterns acquire (i) a surviving working substrate at first exposure,
(ii) postsynaptic recruitment adequate to establish protected
structure, and (iii) episode-distinguishable retrieval --- in that
dependency order, which is exactly the narrowing Phase I measured.
\end{quote}

Three recorded, unexecuted candidate directions belong to this
statement and are listed here only as evidence-adjacent hypotheses,
not as designs: (a) protecting a never-yet-exposed cohort's working
afferents from M2/M4 churn during other blocks (the e1-family question
registered in the source-correction record \cite{fec}; the two gain
registrations supplied the direct evidence that substrate
pre-exposure survival is the binding constraint); (b) a
recruitment mechanism that is not input-proportional at the membrane
(the measured bound \(i_{\mathrm{boost}} \le k_g I_W\) is the reason
both registrations could not engage); and (c) budget partition
re-keyed along the channel-group axis, tested under alternation ---
the only mechanism ever measured to improve coexistence at the
synaptic level, currently tested blocked-only \cite{synmem}. Any of
these, pursued, requires the full discipline: frozen protocol,
identity gate, preserved failures.